\chapter{The Single-Cycle Datapath and Control}

We now build a processor that executes \code{add}, \code{sub}, \code{and},
\code{or}, \code{slt}, \code{lw}, \code{sw}, \code{beq} (and later \code{j})
with each instruction completing in exactly one (long) clock cycle.

\section{Performance}

\begin{definition}
CPU time $=$ Instruction count $\times$ CPI $\times$ Clock cycle time
$= \dfrac{\text{IC} \times \text{CPI}}{\text{Clock rate}}$, where CPI is the
average number of clock cycles per instruction.
\end{definition}

The ``iron law'' shows the trade-offs: the ISA and compiler affect IC, the
implementation affects CPI and cycle time. A single-cycle design has
CPI $= 1$ but a long cycle time.

\section{Datapath Elements}

\begin{itemize}
  \item \textbf{PC} register and an adder computing $PC + 4$.
  \item \textbf{Instruction memory}: address in, instruction out
    (read-only, combinational).
  \item \textbf{Register file}: two read ports, one write port (written at
    the clock edge when \code{RegWrite} = 1).
  \item \textbf{ALU} with \code{Zero} output.
  \item \textbf{Data memory}: address, write data, read data;
    \code{MemRead}/\code{MemWrite} controls.
  \item \textbf{Sign extender} (16 $\to$ 32 bits) and a \textbf{shift-left-2}
    unit for branch offsets, plus a second adder for the branch target.
\end{itemize}

Separate instruction and data memories are needed because a single-cycle
design cannot use one memory twice in the same cycle (a structural hazard).

\section{Building the Datapath}

\paragraph{Fetch.} Instruction $\gets$ IMem[PC]; PC $\gets$ PC + 4
(unless a branch is taken).

\paragraph{R-type.} Read \code{rs}, \code{rt}; ALU computes; write result
to \code{rd}.

\paragraph{Load/store.} Read \code{rs} (base) and, for \code{sw},
\code{rt}; ALU adds base + sign-extended offset; \code{lw} reads data memory
and writes to \code{rt}; \code{sw} writes \code{rt}'s value to memory.

\paragraph{Branch.} ALU subtracts \code{rs} $-$ \code{rt}; if
\code{Zero}, the next PC is $PC + 4 + (\text{signext}(imm) \ll 2)$.

Merging these requires multiplexers wherever a unit has more than one
possible source:
\begin{itemize}
  \item \code{RegDst} mux: write register is \code{rt} (I-type) or
    \code{rd} (R-type).
  \item \code{ALUSrc} mux: second ALU operand is register \code{rt} or the
    sign-extended immediate.
  \item \code{MemtoReg} mux: write-back value is ALU result or memory data.
  \item \code{PCSrc} mux: next PC is $PC+4$ or the branch target;
    $\code{PCSrc} = \code{Branch} \cdot \code{Zero}$.
\end{itemize}

\begin{figure}[H]
\centering
\begin{tikzpicture}[thick, >=Stealth, font=\small,
  blk/.style={draw, minimum height=1.6cm, minimum width=1.4cm, align=center},
  mux/.style={draw, trapezium, trapezium angle=70, shape border rotate=270,
    minimum width=0.9cm, minimum height=0.4cm, inner sep=1pt, font=\scriptsize}]
  \node[blk, minimum width=0.6cm] (pc) at (0,0) {PC};
  \node[blk] (im) at (1.8,0) {Instr\\Mem};
  \node[blk, minimum height=2.2cm] (rf) at (4.4,0) {Reg\\File};
  \node[blk, minimum width=0.9cm] (se) at (4.4,-2.2) {Sign\\ext};
  \node[mux] (asrc) at (6.2,-0.6) {M};
  \node[draw, minimum height=1.8cm, minimum width=1cm] (alu) at (7.4,0) {ALU};
  \node[blk] (dm) at (9.4,0) {Data\\Mem};
  \node[mux] (m2r) at (11,0) {M};
  \node[draw, circle, inner sep=1pt] (add4) at (1.8,2.2) {$+4$};

  \draw[->] (pc) -- (im);
  \draw[->] (im) -- node[above]{\scriptsize rs,rt} (rf);
  \draw[->] (im.south) |- (se.west);
  \draw[->] ($(rf.east)+(0,0.6)$) -- ($(alu.west)+(0,0.6)$);
  \draw[->] ($(rf.east)+(0,-0.6)$) -- (asrc.west);
  \draw[->] (se.east) -| (asrc.south);
  \draw[->] (asrc.east) -- ($(alu.west)+(0,-0.6)$);
  \draw[->] (alu) -- node[above]{\scriptsize addr} (dm);
  \draw[->] (dm) -- (m2r);
  \draw[->] ($(alu.east)+(0,-0.6)$) -- ++(0.3,0) |- ++(0,-1.4) -| (m2r.south);
  \draw[->] (m2r.east) -- ++(0.4,0) |- ++(0,-3.2) -| ($(rf.south)+(0.4,0)$);
  \node[font=\scriptsize] at (8.2,-2.95) {write data};
  \draw[->] (pc.north) |- (add4.west);
  \draw[->] (add4.east) -- ++(1.2,0) node[right]{\scriptsize to PC mux};
\end{tikzpicture}
\caption{Simplified single-cycle datapath (branch logic and RegDst mux
omitted for clarity).}
\end{figure}

\section{Control}

The \textbf{main control unit} decodes the 6-bit opcode into the control
signals. A separate small \textbf{ALU control} unit combines a 2-bit
\code{ALUOp} from main control with the \code{funct} field.

\subsection{Main control table}

\begin{center}
\resizebox{\textwidth}{!}{%
\begin{tabular}{lcccccccc}
\toprule
Instr & RegDst & ALUSrc & MemtoReg & RegWrite & MemRead & MemWrite & Branch & ALUOp \\
\midrule
R-type & 1 & 0 & 0 & 1 & 0 & 0 & 0 & 10 \\
\code{lw} & 0 & 1 & 1 & 1 & 1 & 0 & 0 & 00 \\
\code{sw} & X & 1 & X & 0 & 0 & 1 & 0 & 00 \\
\code{beq} & X & 0 & X & 0 & 0 & 0 & 1 & 01 \\
\bottomrule
\end{tabular}}
\end{center}

X means don't care: \code{sw} and \code{beq} do not write a register, so it
does not matter which register number or write value is selected. (But
\code{RegWrite} and \code{MemWrite} must never be X!)

\subsection{ALU control}

\begin{center}
\begin{tabular}{llll}
\toprule
ALUOp & funct & Operation & ALU control \\
\midrule
00 (lw/sw) & XXXXXX & add & 0010 \\
01 (beq) & XXXXXX & subtract & 0110 \\
10 (R) & 100000 & add & 0010 \\
10 (R) & 100010 & subtract & 0110 \\
10 (R) & 100100 & AND & 0000 \\
10 (R) & 100101 & OR & 0001 \\
10 (R) & 101010 & slt & 0111 \\
\bottomrule
\end{tabular}
\end{center}

This two-level decoding (main control + ALU control) keeps the main
control small and fast.

\subsection{Adding jump}

For \code{j}: the target is $\{(PC+4)[31{:}28], \text{addr}[25{:}0], 00\}$.
Add one more mux after the \code{PCSrc} mux, controlled by a new \code{Jump}
signal from main control.

\begin{example}[Tracing \code{lw \$t1, 8(\$t0)}]
\begin{enumerate}
  \item Fetch from IMem[PC]; adder computes PC+4.
  \item Register file reads \code{\$t0} (rs); \code{rt} is \code{\$t1}.
  \item Sign-extend 8; \code{ALUSrc}=1 selects it; ALU control = add;
    ALU outputs \code{\$t0}+8.
  \item Data memory reads that address (\code{MemRead}=1).
  \item \code{MemtoReg}=1 selects memory data; \code{RegDst}=0 selects
    \code{rt}; \code{RegWrite}=1 writes \code{\$t1} at the clock edge.
  \item \code{Branch}=0, so PC $\gets$ PC+4.
\end{enumerate}
\end{example}

\begin{important}[Exam-style question: control faults]
A common question: ``suppose \code{RegDst} is stuck at 0; which instructions
break?'' Answer: R-type instructions (they would write \code{rt} instead of
\code{rd}); \code{lw} still works; \code{sw}, \code{beq} don't care.
Always reason per instruction using the control table.
\end{important}

\section{Timing of the Single-Cycle Design}

The clock period must accommodate the slowest instruction, which is
\code{lw}:
\[
T = t_{IMem} + t_{RegRead} + t_{ALU} + t_{DMem} + t_{RegWrite\,setup}.
\]

\begin{example}
Let IMem = 200\,ps, register read = 100\,ps, ALU = 200\,ps, DMem = 200\,ps,
register write = 100\,ps (ignore muxes). Instruction times:
R-type 600, \code{lw} 800, \code{sw} 700, \code{beq} 500\,ps. The
single-cycle clock must be 800\,ps even though most instructions need less.
With a mix of 25\% \code{lw}, 10\% \code{sw}, 45\% R-type, 20\% \code{beq},
a hypothetical variable-length clock would average
$0.25(800) + 0.10(700) + 0.45(600) + 0.20(500) = 640$\,ps,
so single-cycle wastes $800/640 = 1.25\times$.
\end{example}

\begin{important}[Problems with single-cycle]
\begin{itemize}
  \item Cycle time set by the slowest instruction; common, fast instructions
    are penalized (violates ``make the common case fast'').
  \item Each functional unit used at most once per instruction, so units
    must be duplicated (two memories, three adders).
  \item Complex instructions (floating point) would make the cycle huge.
\end{itemize}
Solutions: multi-cycle (next chapter) and pipelining.
\end{important}

\section*{Summary}
\begin{itemize}
  \item CPU time = IC $\times$ CPI $\times$ cycle time.
  \item Datapath: PC, IMem, register file, ALU, DMem, sign extend, adders;
    muxes RegDst, ALUSrc, MemtoReg, PCSrc (and Jump).
  \item Control: main control from opcode (table above) + ALU control from
    ALUOp and funct.
  \item Single-cycle: CPI = 1, cycle time = longest instruction (\code{lw}).
\end{itemize}
